\lecture{4}{9. Februar 2026}{Capacitors \& Inductors}

\subsection{Capacitors}

\subsubsection{Capacitance}
\begin{equation*}
  C = \frac{q}{V}
.\end{equation*}

\begin{symbols}
\InsertSym{C}
\InsertSym{q}
\InsertSym{V}
\end{symbols}

\subsubsection{Capacitance of a Parallel-Plate Capacitor}
\begin{equation*}
  C = \epsilon_0 \frac{A}{d}
.\end{equation*}

\begin{symbols}
\InsertSym{C}
\InsertSym{percon}
\InsertSym{Acs}
\InsertSym{d}
\end{symbols}

\subsubsection{I-V Relation for an Ideal Capacitor}
\begin{equation*}
  I_C (t) = C \frac{\mathrm{d}V_C (t)}{\mathrm{d}t} \iff V_C (t) = \frac{1}{C} \int_{-\infty}^{t} I_C (t') \, \mathrm{d}t'
.\end{equation*}

\begin{symbols}
\sym{$I_C$}{Current through capacitor}{\unit{\A}}
\sym{$V_C$}{Voltage over capacitor}{\unit{\V}}
\InsertSym{C}
\InsertSym{t}
\end{symbols}

\subsubsection{Capacitance of a Cylindrical Capacitor}
\begin{equation*}
  C = 2 \pi \epsilon_0 \frac{L}{\ln \frac{b}{a}}
.\end{equation*}

\begin{symbols}
\InsertSym{C}
\InsertSym{percon}
\InsertSym{Llength}
\sym{$a$}{Radius of the inner post}{\unit{\m}}
\sym{$b$}{Radius of the outer post}{\unit{\m}}
\end{symbols}

\subsubsection{Capacitance of a Cylindrical Capacitor}
\begin{equation*}
  C = 4 \pi \epsilon_0 \frac{ab}{b-a}
.\end{equation*}

\begin{symbols}
\InsertSym{C}
\InsertSym{percon}
\sym{$a$}{Radius of the inner sphere}{\unit{\m}}
\sym{$b$}{Radius of the outer sphere}{\unit{\m}}
\end{symbols}


\subsubsection{Parallel Capacitance}
\begin{equation*}
  C_{\mathrm{eq}} = \sum_{i = 1}^{n} C_i
.\end{equation*}

\begin{symbols}
  \sym{$n$}{Number of capacitors}{---}
  \sym{$C_{eq}$}{Equivalent capacitance}{\unit{\F}}
  \sym{$C_i$}{Capacitance of capacitor $i$}{\unit{\F}}
\end{symbols}

\subsubsection{Capacitors in Series}
\begin{equation*}
  \frac{1}{C_{\mathrm{eq}}} = \sum_{i = 1}^{n} \frac{1}{C_i}
.\end{equation*}

\begin{symbols}
  \sym{$n$}{Number of capacitors}{---}
  \sym{$C_{eq}$}{Equivalent capacitance}{\unit{\F}}
  \sym{$C_i$}{Capacitance of capacitor $i$}{\unit{\F}}
\end{symbols}

\subsubsection{Potential Energy in a Capacitor}
\begin{equation*}
  U_E = \frac{1}{2} C V^2
.\end{equation*}

\begin{symbols}
\InsertSym{UE}
\InsertSym{C}
\InsertSym{V}
\end{symbols}

\subsection{Inductors}
\subsubsection{Inductance}
\begin{equation*}
  L = n \frac{\Phi_B}{I}
.\end{equation*}

\begin{symbols}
\InsertSym{Lind}
\InsertSym{nwind}
\InsertSym{PhiB}
\InsertSym{I}
\end{symbols}

\subsubsection{I-V Relation for an Ideal Inductor}
\begin{equation*}
  V_L (t) = L \frac{\mathrm{d}I_L(t)}{\mathrm{d}t} \iff I_L (t) = \frac{1}{L} \int_{t_0}^{t} V_L (t') \, \mathrm{d}t' + I_0
.\end{equation*}


\begin{symbols}
\sym{$I_L$}{Current through inductor}{\unit{\A}}
\sym{$V_L$}{Voltage over inductor}{\unit{\V}}
\InsertSym{Lind}
\InsertSym{t}
\end{symbols}


\subsubsection{Magnetic Field of an Inductor}
\begin{equation*}
  \Vec{B} = \mu_0 I n
.\end{equation*}

\begin{symbols}
\InsertSym{Bvec}
\InsertSym{magper}
\InsertSym{I}
\InsertSym{nwind}
\end{symbols}

\subsubsection{Inductors in Series}
\begin{equation*}
  L_{\mathrm{eq}} = \sum_{i = 1}^{n} L_i
.\end{equation*}

\begin{symbols}
  \sym{$n$}{Number of inductors}{---}
  \sym{$L_{eq}$}{Equivalent Inductance}{\unit{\H}}
  \sym{$L_i$}{Inductance of inductor $i$}{\unit{\H}}
\end{symbols}

\subsubsection{Parallel Inductors}
\begin{equation*}
  \frac{1}{L_{\mathrm{eq}}} = \sum_{i = 1}^{n} \frac{1}{L_i}
.\end{equation*}

\begin{symbols}
  \sym{$n$}{Number of inductors}{---}
  \sym{$L_{eq}$}{Equivalent Inductance}{\unit{\H}}
  \sym{$L_i$}{Inductance of inductor $i$}{\unit{\H}}
\end{symbols}

\subsubsection{Energy Stored in an Inductor}
\begin{equation*}
  U_L = \frac{1}{2} L I^2
.\end{equation*}
\begin{symbols}
\sym{$U_L$}{Potential energy stored in inductor}{\unit{\J}}
\InsertSym{Lind}
\InsertSym{I}
\end{symbols}

\subsubsection{Energy Density in an Inductor}
\begin{equation*}
  u_L = \frac{\Vec{B}}{2 \mu_0}
.\end{equation*}

\begin{symbols}
\sym{$u_L$}{Energy density of inductor}{\unit{\J\per\m^3}}
\InsertSym{Bvec}
\InsertSym{magper}
\end{symbols}

\subsection{RC and RL Circuits in the DC Regime}

\subsubsection{Charge in a (Charging) Capacitor as a Function of Time}
\begin{equation*}
  q(t) = C V_s \left( 1 - e^{- \frac{t}{RC}} \right)
.\end{equation*}

\begin{symbols}
\InsertSym{q}
\InsertSym{C}
\InsertSym{Vs}
\InsertSym{t}
\InsertSym{R}
\InsertSym{C}
\end{symbols}

\subsubsection{Current in a (Charging) Capacitor as a Function of Time}
\begin{equation*}
  I_C(t) = \frac{V_s}{R} e^{- \frac{t}{RC}}
.\end{equation*}

\begin{symbols}
\sym{$I_C$}{Current through capacitor}{\unit{\A}}
\InsertSym{C}
\InsertSym{Vs}
\InsertSym{t}
\InsertSym{R}
\InsertSym{C}
\end{symbols}

\subsubsection{Capacitive Time Constant}
\begin{equation*}
  \tau_C = RC
.\end{equation*}

\begin{symbols}
\InsertSym{tauC}
\InsertSym{R}
\InsertSym{C}
\end{symbols}

\subsubsection{Inductive Time Constant}
\begin{equation*}
  \tau_L = \frac{L}{R}
.\end{equation*}

\begin{symbols}
\InsertSym{tauL}
\InsertSym{Lind}
\InsertSym{R}
\end{symbols}

\subsubsection{Current in an Inductor}
\begin{equation*}
  I = \frac{V_s}{R} \left( 1 - e^{- \frac{t}{\tau_L}} \right)
.\end{equation*}

\begin{symbols}
\InsertSym{I}
\InsertSym{Vs}
\InsertSym{R}
\InsertSym{t}
\InsertSym{tauL}
\end{symbols}
